\chapter{Image spectrale--polaire et réduction exacte de Riemann--Weil}

\section{Point de départ : la structure spectrale--polaire complète}

L'entrée de la résolution n'est pas la forme de Weil isolée ni une liste de
zéros. C'est la structure spectrale--polaire complète produite par le constructeur
commun et attestée par le propriétaire exact qui précède ce chapitre ; ce n'est pas
une donnée externe. La chaîne causale propre à ce chapitre est
\begin{equation}\label{eq:pdf2-rh-cadena-procedencia}
 \mathfrak G_{\rm sp}
 \xrightarrow{\ \operatorname{pr}_{X,{\rm sp}}\ }X_{\chi_{\rm sp}}
 \xrightarrow{\ \mathcal A_{\rm sp}\ }\mathfrak M_{\rm sp}
 \xrightarrow{\ \mathcal P_{\rm sp}\ }\mathcal C_{\rm GW}.
\end{equation}
Le codomaine classique n'apparaît que dans la dernière flèche ; il ne sélectionne pas
rétrospectivement l'état, les cartes, la projection ni le terminal.

La structure de départ contient l'objet à treize coordonnées
\begin{equation}\label{eq:pdf2-rh-res-sp-completa}
\begin{aligned}
 \mathfrak G_{\rm sp}^{(13)}=\bigl(&
 \mathcal F_{\rm sp},
 \operatorname{Ext}_{\Gamma,{\rm sp}},
 \operatorname{Rel}_{\rm sp},
 \mathcal N_{\rm sp},
 \operatorname{or}_{\rm sp},
 \operatorname{Carr}_{\rm sp},
 \operatorname{Mem}_{\rm sp},\\
 &\partial_{\rm sp},
 \gamma_{\rm sp},
 \Sigma_{\rm sp}^{\rm mult},
 \operatorname{Surv}_{\rm sp},
 \tau_{\rm sp}^{\rm term},
 \mathcal L_{\rm sp}\bigr).
\end{aligned}
\end{equation}
Chaque coordonnée intervient matériellement :
\begin{enumerate}
 \item \(\mathcal F_{\rm sp}\) conserve les deux feuillets, leur orientation, la
 paire résidu--quotient et la signature de la publication double.
 \item \(\operatorname{Ext}_{\Gamma,{\rm sp}}\) transporte cette fibre par la
 classe dirigée de coupures et de niveaux nonadiques.
    \item \(\operatorname{Rel}_{\rm sp}\) contient le cocycle de retenue,
    l'orthogonalité \(P=P^*=P^2\), les deux publications et le complément
    structurel \(B_{\rm sp}^{\rm str}:=\Xi^*(I-P)\Xi\). L'identification
    de ce complément à la forme de Guinand--Weil est une flèche distincte
    et fait l'objet d'un audit séparé dans ce chapitre.
 \item \(\mathcal N_{\rm sp}\) enregistre le niveau \(m\), la fenêtre \(R\),
 le quotient \(q=5/9\), les moments des deux feuillets et les indices des
 canaux premier, gamma et polaire.
 \item \(\operatorname{or}_{\rm sp}=\widehat R=(R,K)\) maintient distinctes
 les deux cartes orientées.
 \item \(\operatorname{Carr}_{\rm sp}=\operatorname{Carr}_R\) est la
 retenue opératorielle et ne s'annule pas lors de la publication d'une translation.
 \item \(\operatorname{Mem}_{\rm sp}=(S_0,M,\mathcal O_9)\) conserve
 l'identité finie de Stein et son terminal.
 \item \(\partial_{\rm sp}=(I-P)\Xi\) est la frontière qui n'apparaît pas dans la
 compression \(P\Xi\) et dont l'énergie définit le complément structurel. Son
 identification à la forme de Weil n'est pas incorporée à cette définition.
 \item \(\gamma_{\rm sp}\) est le chemin cofinal de fenêtres, avec un unique
 régulateur pour tous les canaux.
 \item \(\Sigma_{\rm sp}^{\rm mult}\) est la multisection des deux
 publications et de leurs composantes premier--gamma--polaire ; ce n'est pas la signature d'un
 unique état.
 \item \(\operatorname{Surv}_{\rm sp}\) conserve la famille sous
 \(m\mapsto m+1\) et \(\mathcal D_R\hookrightarrow\mathcal D_{R'}\).
 \item \(\tau_{\rm sp}^{\rm term}\) conserve le terminal
 \(M^{*N}S_0M^N\), qui, dans la carte miroir, se publie au moyen de \(E_R\).
    \item \(\mathcal L_{\rm sp}=(\Xi,P,B_{\rm sp}^{\rm str})\) est le lecteur final ; il s'applique
 après antécédent, cartes, mémoire, chemin et terminal.
\end{enumerate}

L'assembleur et le publicateur sont
\begin{equation}\label{eq:pdf2-rh-ensamblador-sp}
\begin{aligned}
 \mathcal A_{\rm sp}(\operatorname{pr}_{X,{\rm sp}}\widehat x)
  ={}&(\widehat X_\infty,\Gamma_9,\widehat R,
      \operatorname{Carr}_R,S_0,M,\mathcal O_9,
      \Xi,P,\mathcal D_R^{\rm cof}),\\
 \mathfrak M_{\rm sp}={}&\operatorname{Im}\mathcal A_{\rm sp},\\
 \mathcal P_{\rm sp}(\mathfrak M_{\rm sp})
  ={}&(B_{\rm sp}^{\rm str},\Xi^*(I-P)\Xi,
       \mathscr I_{\rm GW},\Delta_{\rm bind},
       \text{famille cofinale premier--gamma--polaire}),\\
 \Delta_{\rm bind}:={}&B_{\rm sp}^{\rm str}-\mathscr I_{\rm GW}.
\end{aligned}
\end{equation}
La mémoire finie de l'objet de départ satisfait
\begin{equation}\label{eq:pdf2-rh-stein-finito}
 S_0=\sum_{r=0}^{N-1}M^{*r}\mathcal O_9^*\mathcal O_9M^r
     +M^{*N}S_0M^N.
\end{equation}
Le dernier terme reste visible jusqu'à ce que son extinction soit démontrée. Par
conséquent, remplacer \(\mathfrak G_{\rm sp}\) par l'ombre publiée
\(B_{\rm sp}^{\rm str}\) change l'objet de départ ; ce n'est pas une abréviation
du même théorème. Le résiduel \(\Delta_{\rm bind}\) est conservé jusqu'à ce que la
colligation source--first démontre matériellement son annulation dans
\eqref{eq:pdf2-rh-binding-source-first}.

\section{Domaine cofinal et forme cible}

Pour $R>0$, on fixe
\[
 \mathcal D_R=C_c^\infty((-R/2,R/2)),\qquad
 \widehat\psi(t)=\int_{\mathbb R}\psi(x)e^{-itx}\,dx,
\]
et, pour $\psi,\phi\in\mathcal D_R$,
\[
 h_{\psi,\phi}(z)=\widehat\psi(z)
 \overline{\widehat\phi(\overline z)},\qquad
 \check h_{\psi,\phi}(a)=\langle\psi,T_a\phi\rangle.
\]
L'union $\bigcup_R\mathcal D_R=C_c^\infty(\mathbb R)$ est exacte. Pour tout
ensemble fini de points complexes distincts, les évaluations
$\psi\mapsto(\widehat\psi(z_1),\ldots,\widehat\psi(z_n))$ sont conjointement
surjectives lorsque $R$ est suffisamment grand. Par conséquent, cette classe cofinale
sépare les évaluations requises par le critère de Weil.

\begin{lemma}[Preuve de la séparation conjointe]
\label{lem:pdf2-rh-separacion}
L'application
\[
 \mathcal D_R\longrightarrow\mathbb C^n,
 \qquad
 \psi\longmapsto
 (\widehat\psi(z_1),\ldots,\widehat\psi(z_n))
\]
est surjective pour tout ensemble de points complexes distincts
\(z_1,\ldots,z_n\), dès que \(R\) contient un intervalle non vide.
\end{lemma}

\begin{proof}
Si les fonctionnelles d'évaluation étaient linéairement dépendantes, il existerait
\(c_1,\ldots,c_n\), non tous nuls, tels que
\[
 \int_{\mathbb R}\psi(x)
 \left(\sum_{j=1}^nc_je^{-iz_jx}\right)\,dx=0
 \qquad(\psi\in\mathcal D_R).
\]
La fonction analytique \(\sum_jc_je^{-iz_jx}\) s'annulerait dans
\((-R/2,R/2)\), donc dans tout le plan. L'indépendance des exponentielles
d'exposants distincts impose \(c_j=0\) pour tout \(j\), contradiction.
Les fonctionnelles sont indépendantes ; comme le codomaine est de dimension finie,
l'application est de rang \(\mathbb C^n\).
\end{proof}

\section{Les deux colonnes et les trois canaux calculés avant le signe}

Soit $\vartheta_L$ un régulateur commun. Avec
\[
 \ell_{p,k}=k\log p,\qquad w_{p,k}=(\log p)p^{-k/2},\qquad
 \mu_\Gamma(a)=\frac{e^{-a/2}}{1-e^{-2a}},
\]
on utilise l'inclusion uniforme $j_m$ et le canal de raccordement
$\widehat\delta_{a,m}$, qui satisfont
\[
 j_m^*j_m=I,\qquad
 \widehat\delta_{a,m}^*\widehat\delta_{b,m}
 =(I-T_a)^*(I-T_b).
\]
Les colonnes régularisées sont
\begin{align*}
 J_{+,m,R,L}\psi={}&
 \Bigl(
  (\sqrt{\vartheta_L(\ell_{p,k})w_{p,k}}\,
       \widehat\delta_{\ell_{p,k},m}\psi)_{\ell_{p,k}\leq R},\\
 &\hspace{9mm}a\mapsto
  \sqrt{\vartheta_L(a)\mu_\Gamma(a)}\,
       \widehat\delta_{a,m}\psi,\ \zeta_+b_+(\psi)
 \Bigr),\\
 J_{-,m,R,L}\psi={}&
 \Bigl(
  (\sqrt{2\vartheta_L(\ell_{p,k})w_{p,k}}\,j_m\psi)_{\ell_{p,k}\leq R},
  \sqrt{-c_\Gamma}\,j_m\psi,\ \zeta_-b_-(\psi)
 \Bigr),
\end{align*}
où $c_\Gamma=\psi_\Gamma(1/4)-\log\pi<0$ et
$b_\pm=(\widehat\psi(i/2)\pm\widehat\psi(-i/2))/\sqrt2$.

Un développement terme à terme, sans utiliser de zéros ni le signe recherché, donne
\begin{align*}
 &\langle J_{+,m,R,L}\psi,J_{+,m,R,L}\phi\rangle
 -\langle J_{-,m,R,L}\psi,J_{-,m,R,L}\phi\rangle
 =\mathscr I_{{\rm GW},L}(\psi,\phi),\\
 \mathscr I_{{\rm GW},L}(\psi,\phi)
 ={}&h_{\psi,\phi}(i/2)+h_{\psi,\phi}(-i/2)
 +c_\Gamma\check h_{\psi,\phi}(0)\\
 &+\int_0^\infty\vartheta_L(a)\mu_\Gamma(a)
 [2\check h(0)-\check h(a)-\check h(-a)]\,da\\
 &-\sum_{p,k}\vartheta_L(\ell_{p,k})w_{p,k}
 [\check h(\ell_{p,k})+\check h(-\ell_{p,k})].
\end{align*}
La troncature des deux colonnes en \(\ell_{p,k}\leq R\) est obligatoire, non
notationnelle. Si \(\psi,\phi\in\mathcal D_R\), alors
\(\check h_{\psi,\phi}(\pm\ell)=0\) pour \(\ell>R\) ; la contribution de chaque
paire de coordonnées omise à la \emph{différence} des Grams est donc nulle.
La colonne négative non tronquée aurait, en revanche, une norme infinie car
\(\sum_p(\log p)p^{-1/2}=\infty\). Dans l'ensemble fini
\(\ell_{p,k}\leq R\), les coordonnées premières convergent individuellement lorsque
\(L\to\infty\). Dans le canal gamma, l'intégrande est \(O(a)\) en zéro et
\(\mu_\Gamma(a)=O(e^{-a/2})\) à l'infini, de sorte que la convergence se fait aussi
en norme. Nous définissons donc
\begin{equation}\label{eq:pdf2-rh-columnas-reducidas}
 J_{\pm,m,R}:=\lim_{L\to\infty}J_{\pm,m,R,L}
 \quad\text{dans leurs espaces de Hilbert de coordonnées tronquées},
\end{equation}
et supprimons \(m\) lorsqu'il reste fixe. Par convergence dominée,
\[
 \mathscr I_{{\rm GW},L}\longrightarrow\mathscr I_{\rm GW}
 \qquad(L\to\infty),
\]
et les identités sont naturelles sous $m\mapsto m+1$ et $R\le R'$.

Les colonnes réduites \(J_{+,R}\) et \(J_{-,R}\) sont maintenant des opérateurs
matériels de \(\mathcal D_R\) vers Hilbert. Le développement explicite prouve, sans
consulter de zéros ni de signe,
\begin{equation}\label{eq:pdf2-rh-gram-difference-gw}
 \langle J_{+,R}f,J_{+,R}g\rangle
 -\langle J_{-,R}f,J_{-,R}g\rangle
 =\mathscr I_{\rm pr}(f,g)+\mathscr I_\Gamma(f,g)
  +\mathscr I_{\rm pol}(f,g)
 =\mathscr I_{\rm GW}(f,g).
\end{equation}
Les horloges \(\ell_{p,k}\) produisent le terme premier ; l'intégrale avec
\(\mu_\Gamma\) et \(c_\Gamma\check h(0)\), le terme gamma ; et
\[
 \langle b_+(f),b_+(g)\rangle-
 \langle b_-(f),b_-(g)\rangle
 =h_{f,g}(i/2)+h_{f,g}(-i/2)
\]
produit le terme polaire. Cette partie du calcul est matérielle et indépendante
du pont structurel.

\section{Colligation source--first de l'image spectrale--polaire}

Nous fixons d'abord un niveau nonadique \(m\) ; cet indice est supprimé dans
\(J_{\pm,R}\), \(E_R\) et \(\Sigma_R\) jusqu'à la construction du codeur. La
décomposition des codomaines explicites de
\eqref{eq:pdf2-rh-columnas-reducidas} est
\[
 I_+=\{\mathrm{pr},\Gamma,\mathrm{pol}\},\qquad
 I_-(R)=\{\Gamma0,\mathrm{pol}\}
 \sqcup\{(p,k):\ell_{p,k}\leq R\},
\]
\[
 \mathscr H_{+,R}^{\rm amb}=\bigoplus_{i\in I_+}H_{+,i,R},
 \qquad
 \mathscr H_{-,R}^{\rm amb}=\bigoplus_{j\in I_-(R)}H_{-,j,R}.
\]
Nous n'identifions pas un sous-espace atteignable à une somme de ses projections
coordonnées. Nous définissons, les inclusions ambiantes étant sous-entendues,
\[
 H_{+,R}:=\overline{\operatorname{Ran}J_{+,R}},\qquad
 H_{-,R}:=\overline{\operatorname{Ran}J_{-,R}},
\]
comme sous-espaces fermés des deux Hilbert ambiants. Le propriétaire
spectral imprimé avant ce chapitre construit
\(\mathcal H_R^{\rm sh}\), la différence antispéculaire et la ligne isométrique
source--first au moyen de
\eqref{eq:amp-rh-owner-q}--\eqref{eq:amp-rh-owner-binding}, prouve son
télescopage dans \eqref{eq:amp-rh-owner-telescopia} et publie l'extinction
terminale dans \cref{thm:amp-rh-cierre-cofinal-propietario}. La réalisation
qui suit conserve cette construction, avant l'identification classique,
dans la ligne
\begin{equation}\label{eq:pdf2-rh-coligacion-source-first}
 \Sigma_R=
 \begin{pmatrix}\mathcal K_R^{\rm src}\\ \mathcal F_R^{\rm src}\end{pmatrix}:
 H_{+,R}\longrightarrow
 H_{-,R}\oplus\mathcal H_R^{\rm sh},
 \qquad
 \Sigma_R^*\Sigma_R
 = (\mathcal K_R^{\rm src})^*\mathcal K_R^{\rm src}
  +(\mathcal F_R^{\rm src})^*\mathcal F_R^{\rm src}=I_{H_{+,R}}.
\end{equation}
Son action atteignable est fixée sur la colonne positive complète par
les relations des deux publications qui font partie de
\(\operatorname{Rel}_{\rm sp}\). Pour éviter toute substitution de lecteur,
nous notons
\((\mathcal X_R^{\rm str},\Xi_R^{\rm str},P_R^{\rm str})\) la restriction à
\(\mathcal D_R\) du lecteur structurel de
\eqref{eq:pdf2-rh-res-sp-completa}. Ses deux moments sont, littéralement,
\begin{equation}\label{eq:pdf2-rh-momentos-lector-estructural}
\begin{aligned}
 \langle\Xi_R^{\rm str}f,\Xi_R^{\rm str}g\rangle
  &=\langle J_{+,R}f,J_{+,R}g\rangle,\\
 \langle P_R^{\rm str}\Xi_R^{\rm str}f,
          P_R^{\rm str}\Xi_R^{\rm str}g\rangle
  &=\langle J_{-,R}f,J_{-,R}g\rangle,
 \qquad (P_R^{\rm str})^*= (P_R^{\rm str})^2=P_R^{\rm str}.
\end{aligned}
\end{equation}
La coordonnée de frontière est accompagnée de l'isométrie structurelle
\[
 \jmath_R^{\rm sh}:
 \overline{\operatorname{Ran}((I-P_R^{\rm str})\Xi_R^{\rm str})}
 \longrightarrow\mathcal H_R^{\rm sh}.
\]
Par définition de la ligne source--first, non par une identification ultérieure,
\begin{equation}\label{eq:pdf2-rh-source-state-output}
 E_R:=\mathcal F_R^{\rm src}J_{+,R}
     =\jmath_R^{\rm sh}(I-P_R^{\rm str})\Xi_R^{\rm str}:
 \mathcal D_R\longrightarrow\mathcal H_R^{\rm sh},
 \qquad
 \Omega_{0,R}:=\mathcal K_R^{\rm src}J_{+,R}-J_{-,R}.
\end{equation}
Il n'y a pas une colligation par canal : la même ligne reçoit simultanément les
horloges premier--puissance, l'intégrale gamma, le mode scalaire et les deux lignes
polaires. En particulier, elle conserve tous les produits croisés qui apparaissent en
polarisant \eqref{eq:pdf2-rh-gram-difference-gw}.
Avec les projections ambiantes \(\pi_j^-:\mathscr H_{-,R}^{\rm amb}
\to H_{-,j,R}\), les lignes de sortie correctement typées sont
\begin{equation}\label{eq:pdf2-rh-bloques-coligacion}
 \mathcal K_{j,R}^{\rm src}
 :=\pi_j^-\mathcal K_R^{\rm src}:
 H_{+,R}\longrightarrow H_{-,j,R},
 \qquad
 \mathcal F_R^{\rm src}:H_{+,R}\longrightarrow\mathcal H_R^{\rm sh}.
\end{equation}
Ainsi, on n'applique pas \(\mathcal K_R^{\rm src}\) à une coordonnée positive qui
pourrait ne pas appartenir à elle seule à \(H_{+,R}\). Tous les croisements entre
canaux restent dans
\((\mathcal K_R^{\rm src})^*\mathcal K_R^{\rm src}\) et
\((\mathcal F_R^{\rm src})^*\mathcal F_R^{\rm src}\) ; ils ne sont pas orthogonalisés
avant l'identité d'énergie.

La coordonnée de survie construite par
\eqref{eq:amp-rh-owner-q}, \eqref{eq:amp-rh-owner-telescopia} et
\eqref{eq:amp-rh-owner-terminal} produit les espaces résiduels
\(\mathcal R_{n,R}^{\rm sh}\), les continuations
\(\Omega_{n,R}:\mathcal D_R\to\mathcal R_{n,R}^{\rm sh}\) et les isométries
antispéculaires
\(V_{n,R}:\mathcal R_{n+1,R}^{\rm sh}\to\mathcal R_{n,R}^{\rm sh}\). Son
premier espace et son premier opérateur sont, par définition typée,
\[
 \mathcal R_{0,R}^{\rm sh}:=H_{-,R},\qquad
 \Omega_{0,R}:=\mathcal K_R^{\rm src}J_{+,R}-J_{-,R}.
\]
Sa récurrence d'un tour complet est la publication cofinale du même
opérateur construit, non une donnée ajoutée à la réalisation :
\begin{equation}\label{eq:pdf2-rh-recurrencia-residual}
 q=\frac59,\qquad
 \Omega_{n,R}=qV_{n,R}\Omega_{n+1,R},
 \qquad V_{n,R}^*V_{n,R}=I,
 \qquad
 \sup_{n\geq0}\|\Omega_{n,R}f\|<\infty
 \quad(f\in\mathcal D_R).
\end{equation}
Cette récurrence appartient à la colligation enrichie : elle transporte la
différence de sortie, non la forme de Weil ni son signe.

\begin{proposition}[Annulation coinductive de la différence de sortie]
\label{prop:pdf2-rh-omega-cero}
Pour tout \(R>0\) et tout \(f\in\mathcal D_R\),
\begin{equation}\label{eq:pdf2-rh-output-binding-source}
 \Omega_{n,R}f=0\quad(n\geq0),
 \qquad
 \boxed{\mathcal K_R^{\rm src}J_{+,R}=J_{-,R}}.
\end{equation}
\end{proposition}

\begin{proof}
Itérer \eqref{eq:pdf2-rh-recurrencia-residual} \(N\) fois et utiliser le fait que chaque
\(V_{n,R}\) est isométrique donne
\[
 \|\Omega_{n,R}f\|
 =q^N\|\Omega_{n+N,R}f\|
 \leq q^N\sup_{j\geq n}\|\Omega_{j,R}f\|.
\]
Le membre droit converge vers zéro car \(0<q=5/9<1\). Par conséquent,
\(\Omega_{n,R}f=0\) ; le cas \(n=0\) et
\eqref{eq:pdf2-rh-source-state-output} produisent l'identité encadrée.
\end{proof}

L'action ligne par ligne de la sortie n'est plus implicite. Les projections sont
les restrictions à \(H_{-,R}\) des projections du Hilbert ambiant, et
\begin{equation}\label{eq:pdf2-rh-salidas-generador-a-generador}
\begin{aligned}
 \mathcal K_{\Gamma0,R}^{\rm src}J_{+,R}f
  &=\sqrt{-c_\Gamma}\,j_mf,\\
 \mathcal K_{(p,k),R}^{\rm src}J_{+,R}f
  &=\sqrt{2w_{p,k}}\,j_mf
  &&(\ell_{p,k}\leq R),\\
 \mathcal K_{{\rm pol},R}^{\rm src}J_{+,R}f
  &=\zeta_-b_-(f).
\end{aligned}
\end{equation}
Ces lignes épuisent \(J_{-,R}\), tandis que
\begin{equation}\label{eq:pdf2-rh-frontera-generador-a-generador}
 \mathcal F_R^{\rm src}J_{+,R}f=E_Rf
\end{equation}
conserve, avec tous ses croisements, la coordonnée de frontière.
Évaluer \eqref{eq:pdf2-rh-coligacion-source-first} en \(J_{+,R}f\) et
\(J_{+,R}g\), et utiliser \eqref{eq:pdf2-rh-output-binding-source}, produit
l'identité d'énergie antérieure au signe
\begin{equation}\label{eq:pdf2-rh-energia-source-first}
 \boxed{
 \langle J_{+,R}f,J_{+,R}g\rangle
 =\langle J_{-,R}f,J_{-,R}g\rangle
  +\langle E_Rf,E_Rg\rangle.}
\end{equation}
En la comparant au calcul indépendant
\eqref{eq:pdf2-rh-gram-difference-gw}, on obtient, sans introduire une racine de
la forme cible,
\begin{equation}\label{eq:pdf2-rh-binding-source-first}
 \boxed{E_R^*E_R=\mathscr I_{\rm GW}\big|_{\mathcal D_R}.}
\end{equation}
En particulier,
\(\Delta_{{\rm bind},R}
:=E_R^*E_R-\mathscr I_{\rm GW}|_{\mathcal D_R}=0\) ; cette annulation est la
sortie de \eqref{eq:pdf2-rh-coligacion-source-first}--
\eqref{eq:pdf2-rh-energia-source-first}, non une prémisse insérée dans le
publicateur.

\section{Complément structurel, feuillet miroir et télescopage terminal}

L'état \(E_R=\mathcal F_R^{\rm src}J_{+,R}\) construit par la ligne
source--first est la coordonnée de frontière de
\(\mathfrak G_{\rm sp}\). Nous définissons sa forme structurelle par
\begin{equation}\label{eq:pdf2-rh-bw-estructural}
 \mathcal B_{{\rm sp},R}^{\rm str}
 :=E_R^*E_R
 =\mathscr I_{\rm GW}\big|_{\mathcal D_R},
\end{equation}
où la deuxième égalité est
\eqref{eq:pdf2-rh-binding-source-first}, déjà déduite de la colligation et du
développement générateur par générateur.

Soient \(P_+^{\rm sh}\) et \(P_-^{\rm sh}\) les projecteurs orthogonaux
complémentaires des deux feuillets. On fixe
\begin{equation}\label{eq:pdf2-rh-q-a-d}
 q=\frac59,\qquad
 A=P_+^{\rm sh}+qP_-^{\rm sh},\qquad
 D=\sqrt{1-q^2}\,P_-^{\rm sh}.
\end{equation}
L'orthogonalité donne l'identité exacte
\begin{equation}\label{eq:pdf2-rh-conservacion-local}
 A^*A+D^*D
 =P_+^{\rm sh}+q^2P_-^{\rm sh}
 +(1-q^2)P_-^{\rm sh}=I.
\end{equation}
Les neuf phases coordonnent un seul tour avec mémoire :
\begin{equation}\label{eq:pdf2-rh-monodromia}
 M_9=\Phi A_8\cdots A_0=qV_9,
 \qquad V_9^*V_9=I.
\end{equation}
La contraction par tour est \(q\), non \(q^9\) ; les phases internes conservent
l'ordre et la retenue.

La colligation type
\(E_R:\mathcal D_R\to\operatorname{Ran}P_-^{\rm sh}\). Par conséquent,
\begin{equation}\label{eq:pdf2-rh-binding-espejo}
 P_+^{\rm sh}E_R=0,
 \qquad AE_R=qE_R,
 \qquad
 E_R^*E_R=\mathcal B_{{\rm sp},R}^{\rm str}
 =\mathscr I_{\rm GW}\big|_{\mathcal D_R}.
\end{equation}

\begin{proposition}[Télescopage fini avec terminal conservé]
\label{prop:pdf2-rh-telescopia}
Pour chaque \(N\geq1\),
\begin{equation}\label{eq:pdf2-rh-telescopia-finita}
 \mathcal B_{{\rm sp},R}^{\rm str}
 =\sum_{n=0}^{N-1}(DA^nE_R)^*(DA^nE_R)
 +(A^NE_R)^*(A^NE_R).
\end{equation}
De manière équivalente, évaluée en \(f,g\in\mathcal D_R\), l'identité finie de
la colligation est
\begin{equation}\label{eq:pdf2-rh-balance-finito-columnas}
\boxed{
\begin{aligned}
 &\langle J_{+,R}f,J_{+,R}g\rangle
  -\langle J_{-,R}f,J_{-,R}g\rangle\\
 &\quad=\sum_{n=0}^{N-1}
  \langle DA^nE_Rf,DA^nE_Rg\rangle
  +\langle A^NE_Rf,A^NE_Rg\rangle.
\end{aligned}}
\end{equation}
Le dernier terme satisfait
\begin{equation}\label{eq:pdf2-rh-terminal-extincion}
 (A^NE_R)^*(A^NE_R)
 =q^{2N}\mathcal B_{{\rm sp},R}^{\rm str}
 \xrightarrow[N\to\infty]{}0.
\end{equation}
\end{proposition}

\begin{proof}
En appliquant \eqref{eq:pdf2-rh-conservacion-local} à \(A^nE_R\),
\[
 (A^nE_R)^*(A^nE_R)
 =(DA^nE_R)^*(DA^nE_R)
 +(A^{n+1}E_R)^*(A^{n+1}E_R).
\]
La somme de \(n=0\) à \(N-1\) annule uniquement les termes intermédiaires et
conserve \((A^NE_R)^*(A^NE_R)\). Par
\eqref{eq:pdf2-rh-binding-espejo}, \(A^NE_R=q^NE_R\) ; le terminal
est donc \(q^{2N}E_R^*E_R=q^{2N}\mathcal B_{{\rm sp},R}^{\rm str}\). Comme
\(0<q=5/9<1\), il converge géométriquement vers zéro.
\end{proof}

Ce n'est qu'après avoir démontré \eqref{eq:pdf2-rh-terminal-extincion} que l'on définit
\begin{equation}\label{eq:pdf2-rh-l-r}
 \mathscr L_Rf=(DA^nE_Rf)_{n\geq0}.
\end{equation}
La limite du télescopage fini donne
\begin{equation}\label{eq:pdf2-rh-lstar-l}
 \boxed{
 \mathscr L_R^*\mathscr L_R
 =\mathcal B_{{\rm sp},R}^{\rm str}
 =\mathscr I_{\rm GW}\big|_{\mathcal D_R}\succeq0.}
\end{equation}

\section{Codeur commun, projection et double publication explicites}

Fixons \(m\geq1\) et deux vecteurs orthonormaux
\(u_m,\eta_m\in\mathbb C^{9^m}\), où \(u_m\) est le mode uniforme et
\(\eta_m\) le mode de retenue. L'espace et la projection du codeur sont
\begin{equation}\label{eq:pdf2-rh-hilbert-comun}
 \mathcal X_{m,R}^{\rm sp}
 :=\mathbb C^{9^m}\widehat\otimes
   (H_{-,R}\oplus\mathcal H_R^{\rm sh}),
 \qquad
 P_{m,R}:=|u_m\rangle\langle u_m|\otimes
 I_{H_{-,R}\oplus\mathcal H_R^{\rm sh}},
 \qquad Q_{m,R}:=I-P_{m,R}.
\end{equation}
L'action complète du codeur des cinq canaux est
\begin{equation}\label{eq:pdf2-rh-xi-comun-explicita}
 \boxed{
 \begin{aligned}
  \mathfrak C_{m,R}f
  &:=u_m\otimes(J_{-,R}f,0)
    +\eta_m\otimes(0,E_Rf),\\
  \Xi_{m,R}f&:=\mathfrak C_{m,R}f.
 \end{aligned}}
\end{equation}
On ne prend pas une racine de \(\mathscr I_{\rm GW}\) : les deux entrées sont la
sortie et l'état de la colligation source--first déjà définis dans
\eqref{eq:pdf2-rh-source-state-output}.

\begin{proposition}[Les deux moments s'obtiennent à partir de l'action du codeur]
\label{prop:pdf2-rh-vmas-isometria}
Pour tout \(f,g\in\mathcal D_R\),
\begin{equation}\label{eq:pdf2-rh-dos-momentos-construidos}
\begin{aligned}
 \langle\mathfrak C_{m,R}f,\mathfrak C_{m,R}g\rangle
  &=\langle J_{+,R}f,J_{+,R}g\rangle,\\
 \langle P_{m,R}\mathfrak C_{m,R}f,
         P_{m,R}\mathfrak C_{m,R}g\rangle
  &=\langle J_{-,R}f,J_{-,R}g\rangle,\\
 \langle Q_{m,R}\Xi_{m,R}f,Q_{m,R}\Xi_{m,R}g\rangle
  &=\mathscr I_{\rm GW}(f,g).
\end{aligned}
\end{equation}
\end{proposition}

\begin{proof}
L'orthogonalité \(u_m\perp\eta_m\) et
\eqref{eq:pdf2-rh-energia-source-first} donnent
\[
 \langle\mathfrak C_{m,R}f,\mathfrak C_{m,R}g\rangle
 =\langle J_{-,R}f,J_{-,R}g\rangle
  +\langle E_Rf,E_Rg\rangle
 =\langle J_{+,R}f,J_{+,R}g\rangle.
\]
La projection \(P_{m,R}\) conserve le terme uniforme et élimine celui de
retenue ; \(Q_{m,R}\) fait l'inverse. Les deux autres identités découlent de
\eqref{eq:pdf2-rh-binding-source-first}.
\end{proof}

Le codeur qui vient d'être écrit est un représentant explicite du lecteur
structurel initial, non un lecteur de substitution. Sur les sous-espaces atteignables, on
définit
\begin{equation}\label{eq:pdf2-rh-isomorfismo-lector-estructural}
\begin{aligned}
 \mathcal U_{m,R}^{\rm str}
  (P_R^{\rm str}\Xi_R^{\rm str}f)
   &:=u_m\otimes(J_{-,R}f,0),\\
 \mathcal U_{m,R}^{\rm str}
  ((I-P_R^{\rm str})\Xi_R^{\rm str}f)
   &:=\eta_m\otimes(0,E_Rf).
\end{aligned}
\end{equation}
Les deux sources sont orthogonales. La première égalité de Gram de chaque source
est \eqref{eq:pdf2-rh-momentos-lector-estructural} ; pour la deuxième, on utilise
\(E_R=\jmath_R^{\rm sh}(I-P_R^{\rm str})\Xi_R^{\rm str}\). Ainsi, les deux
règles sont bien définies, sont isométriques et s'étendent à la somme orthogonale
de leurs adhérences. Dans cette adhérence atteignable, elles satisfont
\begin{equation}\label{eq:pdf2-rh-entrelazamiento-lector-estructural}
 \mathcal U_{m,R}^{\rm str}\Xi_R^{\rm str}=\Xi_{m,R},
 \qquad
 \mathcal U_{m,R}^{\rm str}P_R^{\rm str}
 =P_{m,R}\mathcal U_{m,R}^{\rm str},
 \qquad
 (\Xi_R^{\rm str})^*(I-P_R^{\rm str})\Xi_R^{\rm str}=E_R^*E_R.
\end{equation}
Ainsi, \(B_{{\rm sp},R}^{\rm str}\), le complément de la structure de
départ, est exactement le complément du codeur construit ; on n'a pas
changé d'objet en passant de \eqref{eq:pdf2-rh-res-sp-completa} à
\eqref{eq:pdf2-rh-xi-comun-explicita}.

La règle
\begin{equation}\label{eq:pdf2-rh-vmas-rango}
 V_{+,m,R}^{0}(J_{+,R}f):=\Xi_{m,R}f
\end{equation}
est bien définie et isométrique par la première égalité de
\eqref{eq:pdf2-rh-dos-momentos-construidos} ; elle s'étend de manière unique en
une isométrie
\(V_{+,m,R}:H_{+,R}\to\mathcal X_{m,R}^{\rm sp}\). La publication négative
est l'isométrie
\begin{equation}\label{eq:pdf2-rh-vmenos}
 V_{-,m,R}:H_{-,R}\longrightarrow\mathcal X_{m,R}^{\rm sp},
 \qquad V_{-,m,R}y=u_m\otimes(y,0).
\end{equation}
Par construction,
\begin{equation}\label{eq:pdf2-rh-doble-publicacion-material}
 \boxed{
 V_{+,m,R}J_{+,R}=\Xi_{m,R},
 \qquad
 V_{-,m,R}J_{-,R}=P_{m,R}\Xi_{m,R}.}
\end{equation}
Le complément conserve exactement l'état antispéculaire et sa colonne
d'observation :
\begin{equation}\label{eq:pdf2-rh-complemento-comun}
 \boxed{
 \Xi_{m,R}^*Q_{m,R}\Xi_{m,R}
 =E_R^*E_R
 =\mathscr L_R^*\mathscr L_R
 =\mathscr I_{\rm GW}\big|_{\mathcal D_R}.}
\end{equation}

\begin{proposition}[Connecteur induit et coïncidence source--first]
\label{prop:pdf2-rh-conector}
L'opérateur
\begin{equation}\label{eq:pdf2-rh-k-r}
 \mathcal K_R
 :=V_{-,m,R}^*P_{m,R}V_{+,m,R}:
 H_{+,R}\longrightarrow H_{-,R}
\end{equation}
est contractif, coïncide avec \(\mathcal K_R^{\rm src}\) sur tout
\(H_{+,R}\) et satisfait
\begin{equation}\label{eq:pdf2-rh-k-binding}
 \|\mathcal K_R\|\leq1,
 \qquad
 \mathcal K_RJ_{+,R}=J_{-,R}.
\end{equation}
\end{proposition}

\begin{proof}
Les deux cartes sont des isométries et \(P_{m,R}\) est une projection orthogonale,
donc \(\|\mathcal K_R\|\leq1\). Sur l'image dense de \(J_{+,R}\),
\[
 \mathcal K_RJ_{+,R}f
 =V_{-,m,R}^*P_{m,R}\Xi_{m,R}f
 =V_{-,m,R}^*V_{-,m,R}J_{-,R}f
 =J_{-,R}f.
\]
La même identité vaut pour \(\mathcal K_R^{\rm src}\) par
\eqref{eq:pdf2-rh-output-binding-source} ; la continuité et la densité
prouvent que les deux opérateurs coïncident.
\end{proof}

La différence de Gram est ainsi publiée dans toutes ses réalisations
coïncidentes :
\begin{equation}\label{eq:pdf2-rh-tres-publicaciones}
\boxed{
\begin{aligned}
 J_{+,R}^*J_{+,R}-J_{-,R}^*J_{-,R}
 &=\Xi_{m,R}^*Q_{m,R}\Xi_{m,R}\\
 &=E_R^*E_R\\
 &=\mathscr L_R^*\mathscr L_R\\
 &=\mathcal B_{{\rm sp},R}^{\rm str}=\mathscr I_{\rm GW}.
\end{aligned}}
\end{equation}

\section{Naturalité cofinale et critère réversible de Weil}

Pour \(R\leq R'\), soit
\(i_{R,R'}:\mathcal D_R\hookrightarrow\mathcal D_{R'}\). Les colonnes
individuelles ne forment pas un système isométrique : lorsqu'une nouvelle longueur
\(R<\ell_{p,k}\leq R'\) apparaît, les deux acquièrent la même énergie positive. Pour
\(f,g\in\mathcal D_R\), la séparation des supports donne exactement
\begin{equation}\label{eq:pdf2-rh-naturalidad}
 \left\langle(I-T_{\ell_{p,k}})f,
                    (I-T_{\ell_{p,k}})g\right\rangle
 =2\langle f,g\rangle
 =\left\langle\sqrt2j_mf,\sqrt2j_mg\right\rangle .
\end{equation}
Par conséquent, l'incrément premier de
\(J_{+,R'}^*J_{+,R'}\) coïncide avec celui de
\(J_{-,R'}^*J_{-,R'}\) et s'annule dans la différence, mais on ne déclare pas une
fausse isométrie entre les colonnes. La compatibilité cofinale correcte est celle
des formes :
\begin{equation}\label{eq:pdf2-rh-naturalidad-p}
 \boxed{
 \mathscr I_{{\rm GW},R}
 =i_{R,R'}^*\mathscr I_{{\rm GW},R'}i_{R,R'},
 \qquad
 \mathcal B_{{\rm sp},R}^{\rm str}
 =i_{R,R'}^*\mathcal B_{{\rm sp},R'}^{\rm str}i_{R,R'}.}
\end{equation}
Comme \(E_R^*E_R=\mathscr L_R^*\mathscr L_R=\mathscr I_{{\rm GW},R}\), les
réalisations minimales de Kolmogorov admettent bien des isométries uniques sur leurs
images atteignables, définies par
\begin{equation}\label{eq:pdf2-rh-naturalidad-bw}
 \iota_{R,R'}^E(E_Rf):=E_{R'}f,
 \qquad
 \iota_{R,R'}^L(\mathscr L_Rf):=\mathscr L_{R'}f
 \quad(f\in\mathcal D_R).
\end{equation}
Ces deux règles sont bien définies précisément par
\eqref{eq:pdf2-rh-naturalidad-p} ; elles ne s'étendent pas par décret à
\(J_{\pm,R}\), \(\Xi_{m,R}\) ou \(\mathcal K_R\). Le raffinement nonadique
\(m\mapsto m+1\) reste, à fenêtre fixe, dans les isométries structurelles
qui transportent \(j_m\), \(\widehat\delta_{a,m}\), \(u_m\) et \(\eta_m\).

Pour ne pas comprimer la dernière étape, nous fixons aussi ses objets et
normalisations. Soit
\[
 \xi(s):=\frac12s(s-1)\pi^{-s/2}\Gamma(s/2)\zeta(s)
\]
et soit \(\mathscr Z\) le multiensemble de ses zéros, chaque zéro étant répété selon
sa multiplicité. L'équation fonctionnelle et la réalité des coefficients
préservent \(\mathscr Z\) par
\(\rho\mapsto1-\rho\), \(\rho\mapsto\overline\rho\) et
\(\rho\mapsto1-\overline\rho\). Avec notre convention de Fourier, on définit
\begin{equation}\label{eq:pdf2-rh-transformada-weil}
 \Phi_f(s):=\widehat f\!\left(i(s-\tfrac12)\right),
 \qquad
 \widetilde f(x):=\overline{f(-x)}.
\end{equation}
La formule explicite complète calculée dans
\eqref{eq:pdf2-rh-gram-difference-gw}, y compris les termes gamma et polaire
qui retirent les pôles et les zéros triviaux, a la forme spectrale
\begin{equation}\label{eq:pdf2-rh-formula-espectral-weil}
 \mathscr I_{\rm GW}(f,g)
 =\sum_{\rho\in\mathscr Z}
   \Phi_f(\rho)\,
   \overline{\Phi_g(1-\overline\rho)}.
\end{equation}
La somme compte les multiplicités. Elle converge absolument : pour
\(f,g\in C_c^\infty(\mathbb R)\), Paley--Wiener donne une décroissance plus rapide
que toute puissance dans les bandes verticales, tandis que le nombre de zéros avec
\(|\operatorname{Im}\rho|\leq T\) est \(O(T\log T)\).

\begin{proposition}[Critère réversible de Weil avec quantificateurs complets]
\label{prop:pdf2-rh-criterio-weil-desplegado}
Sous la normalisation
\eqref{eq:pdf2-rh-transformada-weil} et en comptant tous les zéros non triviaux
avec leur multiplicité,
\begin{equation}\label{eq:pdf2-rh-criterio-weil}
 \left[
  \mathscr I_{\rm GW}(f,f)\geq0
  \text{ pour tout }f\in C_c^\infty(\mathbb R;\mathbb C)
 \right]
 \Longleftrightarrow
 \left[
  \operatorname{Re}\rho=\frac12
  \text{ pour tout }\rho\in\mathscr Z
 \right].
\end{equation}
\end{proposition}

L'implication réciproque ne s'obtient pas en interpolant un nombre croissant de
zéros : cette opération ne donnerait pas de contrôle uniforme sur la queue. On utilise le
critère classique complet de Weil dans la formulation de Bombieri, dont l'espace
de test et la normalisation sont transposés ici sans les abréger.\footnote{%
E.~Bombieri, \emph{Remarks on Weil's quadratic functional in the theory of
prime numbers, I}, Rendiconti Lincei, Matematica e Applicazioni, ser.~9,
vol.~11 (2000), pp.~183--233. Le critère qui y est fixé affirme que la fonctionnelle
quadratique de la formule explicite est semi-définie positive sur tout
\(C_c^\infty(0,\infty)\) si et seulement si l'hypothèse de Riemann est vraie.}

\begin{proof}
Nous écrivons matériellement le changement de coordonnées qui identifie les deux
critères. Pour \(f\in C_c^\infty(\mathbb R)\), soit
\begin{equation}\label{eq:pdf2-rh-cambio-mellin}
 (\mathcal M f)(x):=x^{-1/2}f(\log x),
 \qquad x>0.
\end{equation}
L’application
\[
 \mathcal M:C_c^\infty(\mathbb R)
 \longrightarrow C_c^\infty(0,\infty)
\]
est bijectif, avec pour inverse
\begin{equation}\label{eq:pdf2-rh-cambio-mellin-inverso}
 (\mathcal M^{-1}g)(u)=e^{u/2}g(e^u).
\end{equation}
Si
\(\widehat g_{\rm Mel}(s):=\int_0^\infty g(x)x^{s-1}\,dx\), le changement
\(x=e^u\) donne, sans facteur résiduel,
\begin{equation}\label{eq:pdf2-rh-mellin-fourier}
 \begin{aligned}
 \widehat{\mathcal M f}_{\rm Mel}(s)
 &=\int_{\mathbb R}f(u)e^{(s-1/2)u}\,du\\
 &=\widehat f\!\left(i(s-\tfrac12)\right)
 =\Phi_f(s).
 \end{aligned}
\end{equation}
De plus,
\begin{equation}\label{eq:pdf2-rh-mellin-involucion}
 \widehat{\overline{\mathcal M f}}_{\rm Mel}(1-\rho)
 =\overline{
   \widehat{\mathcal M f}_{\rm Mel}(1-\overline\rho)}
 =\overline{\Phi_f(1-\overline\rho)}.
\end{equation}
Par conséquent, la fonctionnelle du critère de Weil--Bombieri est exactement
\begin{equation}\label{eq:pdf2-rh-bombieri-identificado}
 \sum_{\rho\in\mathscr Z}
 \widehat{\mathcal M f}_{\rm Mel}(\rho)\,
 \widehat{\overline{\mathcal M f}}_{\rm Mel}(1-\rho)
 =
 \sum_{\rho\in\mathscr Z}
 \Phi_f(\rho)\overline{\Phi_f(1-\overline\rho)}
 =\mathscr I_{\rm GW}(f,f).
\end{equation}
La classe des tests n'est pas modifiée, aucune multiplicité n'est perdue et
l'ensemble des zéros n'est pas tronqué, car \(\mathcal M\) est une bijection. Le
critère classique cité appliqué à
\eqref{eq:pdf2-rh-bombieri-identificado} donne précisément l'équivalence
\eqref{eq:pdf2-rh-criterio-weil}. Dans le sens direct, si
\(\operatorname{Re}\rho=1/2\), la même identité se réduit explicitement à
\[
 \mathscr I_{\rm GW}(f,f)
 =\sum_{\rho\in\mathscr Z}|\Phi_f(\rho)|^2\geq0.
\]
\end{proof}

L'égalité \(\bigcup_R\mathcal D_R=C_c^\infty(\mathbb R)\) est exacte ; chaque
fonction employée dans l'argument appartient à une fenêtre et la naturalité
cofinale y transporte la positivité déjà démontrée.

\begin{theorem}[Résolution de Riemann--Weil à partir de l'image spectrale--polaire]
\label{thm:pdf2-rh-resolucion}
La structure complète
\(\mathfrak G_{\rm sp}\) étant fixée, la forme
complète de Guinand--Weil est positive sur la classe cofinale et séparatrice. En
conséquence, tout zéro non trivial \(\rho\) de la fonction zêta de Riemann
satisfait
\[
 \boxed{\operatorname{Re}\rho=\frac12.}
\]
\end{theorem}

\begin{proof}
Le développement premier--gamma--polaire identifie la différence des deux colonnes
à \(\mathscr I_{\rm GW}\) en utilisant un unique régulateur et le retire par
convergence dominée. La récurrence antispéculaire
\eqref{eq:pdf2-rh-recurrencia-residual} annule la différence de sortie avant de
former une inégalité ; l'isométrie \(\Sigma_R\) donne donc
\eqref{eq:pdf2-rh-energia-source-first}. Comparer cette identité au
développement calculé produit
\(E_R^*E_R=\mathscr I_{\rm GW}\). L'identité locale
\(A^*A+D^*D=I\) est itérée sans effacer le terminal et donne
\eqref{eq:pdf2-rh-telescopia-finita} ; le terminal est exactement
\(q^{2N}E_R^*E_R\) et ne s'éteint qu'à la limite. D'où
\(\mathscr L_R^*\mathscr L_R=\mathscr I_{\rm GW}\succeq0\).
L'action
\eqref{eq:pdf2-rh-xi-comun-explicita} construit alors le codeur, ses
deux moments, la projection et les deux publications, et
\eqref{eq:pdf2-rh-tres-publicaciones} identifie toutes les réalisations du
même carré. La naturalité transporte cette identité à toute l'union
cofinale. Enfin,
\eqref{eq:pdf2-rh-criterio-weil} conclut l'affirmation.
\end{proof}

\section{Contrôle de Redheffer local ultérieur et non directeur}

Le déploiement suivant vérifie que la passivité locale peut être obtenue à partir de
la retenue sans utiliser la forme de Weil. C'est un contrôle de non-circularité, non une
prémisse du \cref{thm:pdf2-rh-resolucion}. Ce transfert tronqué n'est pas
la colligation complète \(\Sigma_R\) de
\eqref{eq:pdf2-rh-coligacion-source-first}.

Pour \(N=9^m\), soit \(U_{a,N}\) l'odomètre unitaire
\[
 U_{a,N}(e_j\otimes f)=e_{j+1}\otimes f\ (j<N-1),
 \qquad U_{a,N}(e_{N-1}\otimes f)=e_0\otimes T_af.
\]
Avec \(P_N=|u_N\rangle\langle u_N|\otimes I\), \(Q_N=I-P_N\),
\(A_{a,N}=P_NU_{a,N}P_N\) et \(C_{a,N}=Q_NU_{a,N}P_N\),
\begin{equation}\label{eq:pdf2-rh-source-energia}
 A_{a,N}^*A_{a,N}+C_{a,N}^*C_{a,N}=P_N,
 \qquad
 C_{a,N}^*C_{a,N}
 =\frac{N-1}{N^2}(2I-T_a-T_a^*).
\end{equation}
Pour \(0<s<1\), le transfert
\begin{equation}\label{eq:pdf2-rh-source-redheffer}
 \mathscr R_{a,N}(s)=(sI-A_{a,N})(I-sA_{a,N})^{-1}
\end{equation}
satisfait
\begin{equation}\label{eq:pdf2-rh-source-defecto}
\begin{aligned}
 I-\mathscr R_{a,N}(s)^*\mathscr R_{a,N}(s)
 ={}&(1-s^2)(I-sA_{a,N}^*)^{-1}\\
 &\quad\cdot C_{a,N}^*C_{a,N}(I-sA_{a,N})^{-1}\succeq0.
\end{aligned}
\end{equation}
Avec
\[
 s_{p,N}=\frac{Np^{-1/2}}{1+(N-1)p^{-1/2}},
 \qquad s_\Gamma(a)=e^{-a/2},
\]
la somme sur les horloges premières et l'intégrale directe gamma conservent une norme au
plus égale à un. Explicitement,
\begin{equation}\label{eq:pdf2-rh-source-relojes}
 \mathscr R_{N,R}
 =\bigoplus_{\log p\leq R}
   \mathscr R_{\log p,N}(s_{p,N})
 \ \oplus\!
 \int_0^\infty{}^\oplus
   \mathscr R_{a,N}(s_\Gamma(a))\,da,
 \qquad \|\mathscr R_{N,R}\|\leq1.
\end{equation}
La carte finie des rôles
\[
 \mathcal H_+=\operatorname{span}\{f_3,f_6,f_9\},\quad
 \mathcal H_-=\operatorname{span}\{f_4,f_8\},\quad
 C_\angle=|f_4\rangle\langle f_3|+|f_8\rangle\langle f_9|
\]
satisfait \(I-C_\angle^*C_\angle=P_6\). Avec les cartes énergétiques
two-way, on fixe
\begin{equation}\label{eq:pdf2-rh-source-cartas}
\begin{aligned}
 \mathcal K_q&=K_2\otimes(\mathbb C^9)^{\otimes q},\\
 \mathsf B_{\pm,q}:\mathcal M_{\pm,q,R}
 &\longrightarrow
 \mathcal K_q\widehat\otimes\mathcal H_\pm
 \widehat\otimes\mathcal H_{\mathfrak C_R},\\
 \mathsf B_{\pm,q}^*\mathsf B_{\pm,q}&=\mathsf G_{\pm,q},
 \qquad
 \mathsf B_{\pm,q}^{\sharp}
 =\mathsf G_{\pm,q}^{-1}\mathsf B_{\pm,q}^*.
\end{aligned}
\end{equation}
Le transfert
\begin{equation}\label{eq:pdf2-rh-source-transferencia}
 \mathscr C_{q,R}^{\rm src}
 =\mathsf B_{-,q}^{\sharp}
  (I\otimes C_\angle\otimes\mathscr R_{N,R})\mathsf B_{+,q}
\end{equation}
a pour défaut
\begin{align}
 \Delta_{q,R}^{\rm src}
 ={}&I\otimes\bigl[P_6\otimes I
 +(P_3+P_9)\otimes
 (I-\mathscr R_{N,R}^*\mathscr R_{N,R})\bigr]\succeq0,
 \label{eq:pdf2-rh-source-delta}\\
 I-(\mathscr C_{q,R}^{\rm src})^\sharp\mathscr C_{q,R}^{\rm src}
 ={}&\mathsf B_{+,q}^{\sharp}\Delta_{q,R}^{\rm src}
 \mathsf B_{+,q}\succeq0.
 \label{eq:pdf2-rh-source-contractividad}
\end{align}
Ce contrôle conserve chemin, horloge, rôle, retenue et canaux avant de former
une différence. Il n'est pas identifié à \(\mathcal K_R\) et ne remplace pas la
publication commune explicite de
\eqref{eq:pdf2-rh-doble-publicacion-material} ; sa portée locale ne
rouvre donc pas la conclusion cofinale.

\section{Réfutateurs d'une ombre amputée}

\begin{falsifier}[Substitutions qui changent l'objet de départ]
L'argument précédent ne se transfère pas à une ombre qui élimine l'une des
treize coordonnées de \eqref{eq:pdf2-rh-res-sp-completa}, construise \(J_+\) et
\(J_-\) à partir d'antécédents distincts, remplace une projection orthogonale par
un idempotent non autoadjoint, omette un canal, utilise des régulateurs incompatibles,
choisisse \(P_R\) après le signe, efface le terminal avant la limite ou réduise
la classe à une famille non cofinale.
\end{falsifier}

\begin{proof}
Chaque opération rompt, respectivement, le domaine de
\(\mathfrak G_{\rm sp}\), la bonne définition de
\eqref{eq:pdf2-rh-vmas-rango}, le carré orthogonal, l'identité fonctionnelle,
la direction causale, le télescopage fini ou le critère réversible. Le
réfutateur affecte l'ombre substituée ; il ne rétrograde pas l'image structurelle
complète employée dans le \cref{thm:pdf2-rh-resolucion}.
\end{proof}

\paragraph{Provenance et force probante.}
La restriction spectrale--polaire, la colligation source--first, son complément et
le feuillet miroir sont \texttt{ARCHITECTURE\_AUTORALE\_PRÉEXISTANTE}.
L'annulation coinductive de \(\Omega\), l'identité finie avec terminal et le
codeur d'action
\eqref{eq:pdf2-rh-xi-comun-explicita} sont
\texttt{FORMALISATION\_NOUVELLE} de cette structure complète. L'identification
Guinand--Weil et la clôture cofinale sont \texttt{RÉSULTAT\_RÉCUPÉRÉ}. Les
équations
\eqref{eq:pdf2-rh-coligacion-source-first}--
\eqref{eq:pdf2-rh-tres-publicaciones} exhibent l'objet, ses domaines, son
action, le terminal et la limite sans introduire une racine abstraite de la forme ;
leur force est \emph{démontré avec structure de départ explicite}.
